\documentclass[tikz,border=7pt]{standalone}
\usepackage[T1]{fontenc}
\usepackage[utf8]{inputenc}
\usepackage[italian]{babel}
\usepackage{tikz}
\usetikzlibrary{arrows.meta,positioning}

\begin{document}
\begin{tikzpicture}[
    level distance=16mm,
    level 1/.style={sibling distance=47mm},
    level 2/.style={sibling distance=26mm},
    level 3/.style={sibling distance=18mm},
    level 4/.style={sibling distance=14mm},
    every node/.style={font=\scriptsize, align=center},
    trie node/.style={circle, draw, fill=black!5, minimum size=8mm, inner sep=1pt},
    new node/.style={circle, draw=orange!80!black, dashed, fill=orange!10, minimum size=8mm, inner sep=1pt},
    edge from parent/.style={draw, line width=0.45pt},
    edge label/.style={font=\scriptsize, inner sep=1pt, fill=white},
    path edge/.style={draw=blue!65!black, line width=1.1pt},
    note/.style={rectangle, draw=black!30, rounded corners=2pt, fill=black!3, align=left, font=\scriptsize, inner sep=5pt}
]

% Trie costruito dal dizionario dell'esempio LZ78 su abaababaa.
\node[trie node] (n0) {$0$\\$\epsilon$}
    child {
        node[trie node] (n1) {$1$\\\texttt{a}}
        child {
            node[trie node] (n3) {$3$\\\texttt{aa}}
            edge from parent node[edge label, left] {\texttt{a}}
        }
        edge from parent node[edge label, left] {\texttt{a}}
    }
    child {
        node[trie node] (n2) {$2$\\\texttt{b}}
        child {
            node[trie node] (n4) {$4$\\\texttt{ba}}
            child {
                node[trie node] (n5) {$5$\\\texttt{baa}}
                child {
                    node[new node] (n6) {$6$\\\texttt{baab}}
                    edge from parent node[edge label, right] {\texttt{b}}
                }
                edge from parent node[edge label, left] {\texttt{a}}
            }
            edge from parent node[edge label, right] {\texttt{a}}
        }
        edge from parent node[edge label, right] {\texttt{b}}
    };

% Evidenzia il cammino seguito codificando baab.
\draw[path edge] (n0) -- (n2);
\draw[path edge] (n2) -- (n4);
\draw[path edge] (n4) -- (n5);
\draw[-Latex, orange!80!black, line width=0.9pt] (n5) -- (n6);

\node[note, right=18mm of n4] {
    input \texttt{baab}\\
    ultimo nodo noto: $5=\texttt{baa}$\\
    prossimo carattere: \texttt{b}\\
    output: $\langle 5,\texttt{b}\rangle$
};

\end{tikzpicture}
\end{document}
